% =====================================================================
% 1. Methodology and experimental design
% =====================================================================
\section{Methodology and experimental design}\label{sec:method}

\subsection{From the guidelines to a plan}

The professor's guidelines ask for three areas (core architectures, learning strategies, basic
generative AI), mandatory ablations, simple models and methodological clarity
\cite{course_message_2011}. We kept one use case for the three areas: the five-superclass task of
PTB-XL \cite{wagner2020ptbxl}. One model serves two areas: the autoregressive generator of lead II
(area 3) is also the model pre-trained without labels and transferred to the classification task
(area 2). The project was planned as a list of tasks with an owner and a reviewer
(\texttt{docs/PLAN\_DL.md}); each block was decided before it was coded (\texttt{docs/DECISIONS.md},
entries D0--D20) and ended with a block note (\texttt{docs/blocks/}).

\subsection{Protocol shared by every block}

\begin{itemize}[leftmargin=1.2em,itemsep=1pt]
  \item \textbf{Same data, loss, metric and logging for every model}, so that differences come from the
        model. Official folds: training 1--8, validation 9, test 10; records without a superclass left
        out; per-lead z-score with training statistics.
  \item \textbf{Baseline and one factor at a time.} Each ablation changes one thing with respect to a
        named reference (the baseline, or another configuration when the question is an interaction, e.g.\
        ``6 blocks without residual'' against ``6 blocks'').
  \item \textbf{Three seeds for every configuration} (42, 43, 44). The first plan was one seed for ablations
        and three for final models; from the Dense block on we decided three for every configuration, so
        that each ablation is compared with its reference seed by seed.
  \item \textbf{Equal, small budgets.} Adam, learning rate $10^{-3}$, batches of 64, at most 50 epochs,
        early stopping with patience 5 on the validation selection metric; gradient clipping at norm 1
        for recurrent, attention and generative models. Exceptions are written in the decisions log with
        their reason (e.g.\ batches of 32 records in the transfer experiments, for GPU memory).
  \item \textbf{Selection on validation, test once.} The rule that picks the final model of a block was
        written before its runs; the test fold is read only by a \texttt{--final} run, and the training
        code refuses a second final run of the same model, tag and seed.
  \item \textbf{Paired comparisons.} Every run stores its predictions, so any two runs on the same split
        can be compared with the paired patient bootstrap (Section~\ref{sec:training}).
\end{itemize}

\subsection{Working practices that saved time}

\begin{itemize}[leftmargin=1.2em,itemsep=1pt]
  \item \textbf{Smoke mode.} \texttt{--fast} trains on 512 records for 2 epochs with 50 bootstrap
        resamples. It runs the whole pipeline of a block in minutes, writes outputs whose names end in
        \texttt{\_fast} (ignored by git) and is refused by the script that generates the numbers.
  \item \textbf{Estimate before long runs.} The longer blocks have an \texttt{estimate} command that
        measures time per batch and peak GPU memory of every configuration before the grid is launched.
        It prevented memory errors in training, but not in generation, which it did not cover
        (Appendix~\ref{app:ai-code}, Sheet C1).
  \item \textbf{Determinism.} Python, NumPy and PyTorch are seeded before the model is built, the same
        seed orders the batches, cuDNN runs in deterministic mode, and recurrent and attention models set
        \texttt{CUBLAS\_WORKSPACE\_CONFIG} as the PyTorch documentation of \texttt{nn.LSTM} requires.
        In the RNN and Transformer blocks the final runs reproduced their ablation runs exactly (same
        validation AUROC and best epoch).
  \item \textbf{Runs, then report.} Each block has \texttt{run} (training only, skipping what was already
        run), \texttt{final} and \texttt{report} (summary log, tables, figures), so a figure can be
        redrawn without retraining.
\end{itemize}

\subsection{What the project trained}

\begin{table}[htbp]
\centering
\caption{Runs of the project, per block (all with seeds 42--44).}
\label{tab:runs}
\small
\begin{tabularx}{\textwidth}{>{\raggedright\arraybackslash}p{4.2cm}>{\raggedright\arraybackslash}X>{\raggedright\arraybackslash}p{2.7cm}}
\toprule
Block (owner) & Configurations & Test fold \\
\midrule
Dense (OF) & baseline + 10 ablations; final \texttt{no\_norm} & final model only \\
1D CNN (OF) & baseline + 9 ablations; final \texttt{no\_dropout} & final model only \\
RNN and RNN + attention (MM) & baseline + 11 ablations; finals \texttt{conv} and \texttt{raw\_attn} & final models only \\
Transformer (MM) & baseline + 14 configurations (step by step and ablations); final \texttt{no\_pos} & final model only \\
Generators (OF) & ECG-GPT baseline + 7 ablations & never read \\
Transfer (OF + MM) & 3 strategies $\times$ 2 label fractions; 2 learning-rate ablations $\times$ 2 fractions; probes on 6 other generators $\times$ 2 fractions & 6 main configurations \\
12 leads vs lead II (OF + MM) & final model of each of the 5 families on lead II & every run \\
Zero-shot (OF + MM) & CLIP, no training & one evaluation \\
\bottomrule
\end{tabularx}
\end{table}
